\section{PHASE I: COARSE IDENTIFICATION}
\label{sec:phase1}

The first phase aims to solve a difficult task with a simple tool: uniform random exploration. The goal is for every player to obtain a coarse but reliable confidence bracket on each cell's maximum value, $\mu^\ast(C)$. The exploration is ``collision-censored"-players make no attempt to avoid each other and instead rely on randomness to provide a sufficient number of non-collision events to learn from.

\subsection{Phase I Protocol}




This phase runs for a fixed budget of $T_0$ rounds. In each round $t \in [T_0]$, every player $j$ independently samples a cell $C_t^{(j)}$ uniformly at random from $\Pcal$ and probes its center $x_{C_t^{(j)}}$ and observes a reward $Y_t^{(j)}$ if she is the unique occupant of the cell, or the null symbol $\bot$ otherwise.

Let $U_j(C,t)$ be the indicator that player $j$ uniquely occupies cell $C$ at round $t$. We track the \emph{success count} for each player-cell pair:
\begin{equation}
\label{eq:success-count}
o_{j,C}(T_0) := \sum\nolimits_{t=1}^{T_0} U_j(C,t).
\end{equation}
If $U_j(C,t)=1$, the player observes a reward $Y_t^{(j)}$; otherwise, she receives the null symbol $\bot$ . The per-round success probability $p_K$ remains constant. Since players sample independently, $p_K$ is the probability that player $j$ selects $C$ while all others avoid it:
\begin{equation}
\label{eq:pk-def}
p_K := \tfrac{1}{K}\left(1-\frac{1}{K}\right)^{N-1}.
\end{equation}




From these successful observations, each player computes an empirical mean for each cell's center, $\widehat{\mu}_{j}^{(0)}(C)$. This estimate, combined with a concentration radius $r^{(0)}_{j}(C)$ and the geometric bracket from the Lipschitz property, yields initial lower and upper confidence bounds on the cell's true maximum value, $\mu^\ast(C)$. 

Formally, the initial confidence brackets for each cell maximum $\mu^*(C)$ are:
\begin{equation}
\label{eq:phase1-brackets}
\begin{aligned}
\LCB^{(0)}_j(C) &:= \widehat{\mu}_{j}^{(0)}(C) - r^{(0)}_{j}(C), \\
\UCB^{(0)}_j(C) &:= \widehat{\mu}_{j}^{(0)}(C) + r^{(0)}_{j}(C) + (Lh\sqrt{d})/2,
\end{aligned}
\end{equation}
where $r^{(0)}_{j}(C)$ is the concentration radius defined with $\beta_0 := \log \frac{4NK(T_0+1)}{\delta_I}$:
\begin{equation}
\label{eq:r0}
r^{(0)}_{j}(C) := \sqrt{\frac{\beta_0}{2\,\max\{1, o_{j,C}(T_0)\}}}.
\end{equation}


\subsection{Phase I Guarantees}
Despite the collision-prone nature of the exploration, standard concentration inequalities show that this simple protocol is highly effective. The first result establishes that the number of successes, $o_{j,C}(T_0)$, is sharply concentrated around its mean for all players and cells simultaneously. The second shows that the empirical means are accurate estimates of the true center values.

\begin{lemma}[Success Counts Under Collisions]
\label{lem:counts}
For any $\eta\in(0,1)$, with probability at least $1-\delta_{I}/2$, the success count for every player $j\in[N]$ and every cell $C\in \Pcal$ is bounded by $(1\pm\eta) T_0 p_K$.
\end{lemma}

\begin{lemma}[Anytime Concentration for Center Means]
\label{lem:means}
With probability at least $1-\delta_{I}/2$, for every player $j\in[N]$ and every cell $C\in \Pcal$, the empirical mean is close to the true mean: $|\widehat{\mu}_j^{(0)}(C) - \mu(x_C)| \le r^{(0)}_{j}(C)$. This holds for any realized value of the success count $o_{j,C}(T_0)$.
\end{lemma}

The proofs of these lemmas, which involve standard applications of Chernoff and Hoeffding bounds with a union bound over all players and cells, are deferred to the appendix. By combining these statistical guarantees with the geometric bound derived from the Lipschitz property, we arrive at the main result of Phase I: with high probability, every player constructs a valid confidence interval for every cell's true maximum value.

\begin{proposition}[Phase-I Maxima Brackets]
\label{prop:phase1-brackets}
With probability at least $1-\delta_{I}$, for every player $j\in[N]$ and every cell $C\in \Pcal$, the computed bounds are valid: $\LCB^{(0)}_j(C) \le \mu^\ast(C) \le \UCB^{(0)}_j(C)$.
\end{proposition}
